\documentclass[11pt]{article}

\usepackage[margin=1in]{geometry}
\usepackage{amsfonts,amsmath}
\usepackage{booktabs}
\usepackage{multirow}
\usepackage{graphicx}
\usepackage{xcolor}
\usepackage[colorlinks,citecolor=blue,urlcolor=blue,linkcolor=blue]{hyperref}
\usepackage[numbers,sort&compress]{natbib}
\usepackage[affil-it]{authblk}
\usepackage{lineno}

\renewcommand{\arraystretch}{1.3}
\emergencystretch=1.5em

\newcommand{\Gr}{\mathrm{Gr}}
\newcommand{\R}{\mathbb{R}}

\title{Ranking Without Certification: Standard Single-Cell Embedding Metrics\\
Cannot Detect Biological Signal at Foundation-Model Dimensionality}

\author{Elliot Tower}
\affil{Independent Researcher \\ \texttt{elliot@elliottower.ai}}
\date{}

\begin{document}
\linenumbers
\maketitle

\begin{abstract}
Single-cell embedding metrics are used for two jobs the field conflates:
\emph{ranking} models against each other, and \emph{certifying} that a
model captures biological signal at all. We show these come apart. CKA
and Procrustes on cell-type centroids rank embeddings by downstream
transfer (partial $\rho = +0.53$ and $+0.66$), but their expected value
under a random Gaussian null follows $\mathbb{E}[\mathrm{CKA}] \approx
1 - 1.06\,(k{-}1)/(d{+}k)$, placing a computable ceiling above all
achievable scores at foundation-model dimensionality ($d \geq 512$). To
push the null below the observed CKA mean of 0.74 at $d = 512$ requires
$k > 168$ shared cell types---beyond any existing dataset. The boundary
between regimes where CKA can and cannot certify falls inside the
dimensionality range the field uses: at $d = 50$, 86\% of conditions
across four datasets exceed the null; at $d = 200$, 17\%; at
$d \geq 512$, none. A PCA-reduction control confirms dimensionality is
causal: Geneformer embeddings PCA-reduced from $d = 512$ to $d = 50$
clear the null in 18/22 tissues, versus 0/22 at native dimensionality.
Two attempted repairs fail in diagnosable ways: a pre-registered
spectral-gap statistic inverts (baselines more concentrated than trained
models), and $z$-score normalization degrades ranking by removing real
$k$-dependent signal at low $d$. Bio-conservation metrics show
non-monotonic response to controlled degradation (single-seed;
multi-seed replication planned). A cross-tissue replication experiment
localizes the ranking failure: scIB metrics achieve 22\% pairwise
inversions on cross-tissue transfer ($\rho = +0.63$, $p = 0.0014$,
stable across three seeds), compared with 30--48\% inversions in the
cross-assay setting where protocol artifacts confound. The failure is
protocol-dependent, not intrinsic. We propose a usability criterion:
given $k$ shared cell types and dimension $d$, compute
$\mathbb{E}[\mathrm{CKA}]$; if scores do not exceed it, CKA is
uninformative as a certification tool. Every current foundation model
operates in the regime where this criterion is violated.
\end{abstract}

\noindent\textbf{Keywords:} single-cell foundation models, embedding evaluation, construct validity, centered kernel alignment, Procrustes analysis, scIB

\section{Introduction}
\label{sec:intro}

Single-cell foundation model benchmarks consistently find that PCA on
raw counts matches or exceeds large pretrained models across most
evaluation axes \citep{wu2025scfm, vcbench2026, unifiedscfm2026,
cellxgeneintegration2026}. \citet{wu2025scfm} raised the possibility
that the metrics themselves might be at fault. The question is whether
high metric scores reflect biological signal learned by the model, or
geometric properties that arise in any high-dimensional embedding.

Embedding metrics serve two roles that the field treats as one. The
first is \emph{ranking}: given two trained models, which produces a
better embedding? The second is \emph{certification}: does a given
model capture biological signal at all? These are logically independent.
A metric can rank trained models correctly while scoring all of them
below unstructured random noise, or it can separate trained from random
while failing to rank among trained models. Treating one as evidence for
the other conflates two distinct claims about metric quality.

We evaluate twelve metrics spanning three families---scIB
bio-conservation \citep{luecken2022scib}, cell-type centroid alignment,
and five geometric modules---on 104 embedding conditions across 25
tissues. Ground truth is cross-technology cell-type transfer F1 from
10x Chromium 3$'$ v3 to Smart-seq2.

Five findings emerge. First, most metrics rank well: Procrustes
similarity on cell-type centroids achieves partial $\rho = +0.66$ with
downstream transfer, and nine of twelve metrics reach significance after
controlling for dimensionality. Second, the CKA certification ceiling is
a closed-form function of $k$ and $d$:
$\mathbb{E}[\mathrm{CKA}] \approx 1 - 1.06\,(k{-}1)/(d{+}k)$. This
formula places the boundary between regimes where CKA can and cannot
certify: at $d = 50$, 86\% of conditions across four datasets clear the
null; at $d \geq 512$, none do. A PCA-reduction control confirms
dimensionality is causal independent of model architecture. Third, two
attempted repairs fail in diagnosable ways: a pre-registered
spectral-gap statistic inverts (null baselines more concentrated than
trained models), and $z$-score normalization degrades ranking by
removing real $k$-dependent signal. Fourth, bio-conservation metrics
show non-monotonic response to controlled embedding degradation, with
scores increasing at intermediate noise levels before collapsing
(single-seed; multi-seed replication planned). Fifth, a cross-tissue
replication experiment localizes the ranking failure: scIB metrics
achieve 22\% pairwise inversions on cross-tissue transfer
($\rho = +0.63$, $p = 0.0014$), compared with 30--48\% in the
cross-assay setting where protocol artifacts confound.

These findings yield a usability criterion: given $k$ shared cell types
and dimension $d$, compute the analytic null; if the expected effect
size does not exceed it, CKA is uninformative as a certification tool.
Every current single-cell foundation model operates in the regime where
this criterion is violated, while the PCA-based methods the field used
previously were on the correct side of the boundary.


\section{Results}

\subsection{Evaluation panel}
\label{sec:panel}

We constructed a cross-technology transfer panel from CellxGene Census
(v2023-12-15). Source embeddings are computed on 10x Chromium 3$'$ v3
cells; target labels come from Smart-seq2 (EFO:0008931) cells in the
same tissue with zero donor overlap. Tissues were auto-discovered by
requiring $\geq 200$ cells per assay and $\geq 8$ shared cell types,
yielding 25 tissues. Eight embedding models
span three architectural families: Geneformer \citep{theodoris2023geneformer}
($d = 512$), Geneformer-v2 104M and 316M, scGPT \citep{cui2024scgpt}
($d = 512$), scVI \citep{lopez2018scvi} ($d = 50$), bag-of-genes
PCA-512 ($d = 512$; following \citealp{cole2026morpheus}), random
Gaussian projection ($d = 512$), and an untrained two-layer MLP encoder
($d = 512$). Six contender models (excluding the two null baselines)
produce 104 contender conditions. Ground truth is macro-averaged
logistic-regression F1 on shared cell types, trained on source
embeddings and evaluated on target.


\subsection{Metric prediction of cross-technology transfer}
\label{sec:prediction}

Twelve metrics were evaluated on all 104 contender conditions. Partial
Spearman rank correlations control for embedding dimensionality via rank
residuals from OLS; tissue-stratified permutation tests (10,000
iterations) provide $p$-values robust to within-tissue dependence.

Nine of twelve metrics significantly predict cross-technology transfer
(Table~\ref{tab:metrics}). Procrustes similarity on cell-type centroids
achieves the strongest correlation (partial $\rho = +0.66$,
$p < 10^{-13}$), followed by target-domain silhouette ($\rho = +0.65$)
and silhouette on cell-type labels ($\rho = +0.60$). Cell-type CKA
achieves $\rho = +0.53$ ($p < 10^{-8}$). Among scIB metrics, graph
connectivity ($\rho = +0.58$), NMI ($\rho = +0.55$), and kNN purity
($\rho = +0.52$) predict transfer at comparable strength. ARI predicts
more weakly ($\rho = +0.37$). The three failures are
direction stability ($\rho = +0.02$, $p = 0.83$) and
participation ratio ($\rho = +0.45$, nominally significant but
confounded by dimensionality; Supplementary~\ref{app:modules}).

\begin{table}[t]
\centering
\caption{Partial Spearman correlation of twelve embedding metrics with
cross-technology transfer F1 ($n = 104$ contender conditions, 25
tissues). Controlling for embedding dimensionality $d$; $p$-values from
tissue-stratified permutation (10,000 iterations). Metrics ranked by
$|\rho|$. Bold: $p < 0.001$.}
\label{tab:metrics}
\small
\begin{tabular}{@{}lccl@{}}
\toprule
Metric & Partial $\rho$ & Perm.\ $p$ & Family \\
\midrule
Procrustes similarity   & $+0.652$ & $\mathbf{< 10^{-4}}$ & centroid \\
Silhouette (target)     & $+0.647$ & $\mathbf{< 10^{-4}}$ & scIB bio \\
Silhouette (label)      & $+0.599$ & $\mathbf{< 10^{-4}}$ & scIB bio \\
Graph connectivity      & $+0.579$ & $\mathbf{< 10^{-4}}$ & scIB bio \\
NMI (Leiden)            & $+0.553$ & $\mathbf{< 10^{-4}}$ & scIB bio \\
Cell-type CKA           & $+0.530$ & $\mathbf{< 10^{-4}}$ & centroid \\
kNN purity              & $+0.515$ & $\mathbf{< 10^{-4}}$ & scIB bio \\
cLISI                   & $+0.493$ & $\mathbf{< 10^{-4}}$ & scIB bio \\
Participation ratio     & $+0.450$ & $\mathbf{< 10^{-4}}$ & geometric \\
Isolated label ASW      & $+0.425$ & $\mathbf{< 10^{-4}}$ & scIB bio \\
ARI (Leiden)            & $+0.368$ & $\mathbf{< 10^{-4}}$ & scIB bio \\
Direction stability     & $+0.022$ & $0.83$               & geometric \\
\bottomrule
\end{tabular}
\end{table}

Controlling for both dimensionality and the number of shared cell types
attenuates correlations by $\leq 3\%$ (CKA: $\rho = 0.52$; Procrustes:
$\rho = 0.64$), confirming that the predictions are driven by embedding
quality rather than panel structure.


\subsection{The CKA null floor is a function of dimensionality}
\label{sec:nullfloor}

Predictive validity does not imply discriminative validity. We
constructed a null baseline by generating pairs of random Gaussian
centroid matrices ($k \times d$, entries $\sim \mathcal{N}(0,1)$) and
computing CKA and Procrustes similarity. The expected CKA under this
null admits a closed-form expression via Wishart moments
(Methods~\ref{sec:closedform}):
\begin{equation}
\label{eq:cka_null}
\mathbb{E}[\mathrm{CKA}(X, Y)] \approx 1 - 1.06 \cdot \frac{k-1}{d+k}
\end{equation}
This formula matches simulation within 0.02\% across $k = 8$--$44$
(Table~\ref{tab:nullfloor}). The formula depends only on $k$ and $d$,
so the null floor is determined by the embedding dimension and the
number of cell types, independent of the model, tissue, or biology.

The formula reveals that the null floor is not a constant: it collapses
as $d$ decreases. At $d = 512$, $k = 16$, the null is 0.970; at
$d = 50$, $k = 16$, it drops to 0.772. To push the null below the
observed CKA mean of 0.74 requires $k > 168$ shared cell types at
$d = 512$, $k > 66$ at $d = 200$, and only $k > 18$ at $d = 50$
(Table~\ref{tab:crossover}). No single-cell dataset has 168 shared cell
types across technology pairs, so CKA at $d = 512$ is confounded by
construction.

\begin{table}[t]
\centering
\caption{The CKA null floor as a function of $d$ and $k$. Analytic null
$\mathbb{E}[\mathrm{CKA}] \approx 1 - 1.06\,(k{-}1)/(d{+}k)$,
verified by simulation (1,000 draws per cell). The null floor tightens
monotonically with $d/k$; at foundation-model dimensionalities
($d \geq 512$), it exceeds all observed embedding CKA scores.}
\label{tab:nullfloor}
\small
\begin{tabular}{@{}rcccc@{}}
\toprule
& \multicolumn{4}{c}{Analytic CKA null at dimension $d$} \\
\cmidrule(lr){2-5}
$k$ & $d = 50$ & $d = 200$ & $d = 512$ & $d = 1280$ \\
\midrule
8  & 0.872 & 0.964 & 0.986 & 0.994 \\
14 & 0.785 & 0.936 & 0.974 & 0.989 \\
22 & 0.691 & 0.900 & 0.958 & 0.983 \\
44 & 0.515 & 0.813 & 0.921 & 0.965 \\
\bottomrule
\end{tabular}
\end{table}

\begin{table}[t]
\centering
\caption{Crossover: number of shared cell types $k$ needed for the
analytic null to drop below the observed CKA mean of 0.74, as a function
of embedding dimension $d$.}
\label{tab:crossover}
\small
\begin{tabular}{@{}rrl@{}}
\toprule
$d$ & $k$ needed & Feasibility \\
\midrule
50   & $> 18$  & Routine \\
200  & $> 66$  & Rare but achievable \\
512  & $> 168$ & Beyond any existing dataset \\
1280 & $> 417$ & Infeasible \\
\bottomrule
\end{tabular}
\end{table}


\subsection{Discriminability across dimensionality regimes}
\label{sec:deffect}

The analytic null generates a testable prediction: conditions at low $d$
should clear the null floor, while conditions at high $d$ should not. We
tested this across four datasets spanning $d = 50$ to $d = 1{,}152$
(Table~\ref{tab:deffect}). At $d = 50$, 218 of 253 conditions (86\%)
exceed their $(k, d)$-matched null. At $d = 200$, 76 of 456 (17\%).
At $d \geq 512$, zero of 129. The gradient is monotone across four
levels of $d$ and holds within each dataset independently.

\begin{table}[t]
\centering
\caption{Fraction of conditions exceeding the analytic CKA null, by
embedding dimension. Four datasets: Census (25 tissues, 8 models),
pancreas (56 pairs), Tabula Sapiens (16 pairs), PBMC (42 pairs).}
\label{tab:deffect}
\small
\begin{tabular}{@{}rrrrl@{}}
\toprule
$d$ & Above null & Total & \% above & Source \\
\midrule
50   & 218 & 253 & 86.2 & Census (scVI) + external PCA \\
200  &  76 & 456 & 16.7 & External PCA/random \\
512  &   0 & 125 &  0.0 & Census (all $d{=}512$ models) \\
$\geq 768$ & 0 & 4 & 0.0 & Census (Geneformer-v2) \\
\bottomrule
\end{tabular}
\end{table}

This gradient could reflect dimensionality or model family, since the
$d = 50$ conditions include scVI and PCA while the $d = 512$ conditions
include foundation models and PCA. A within-panel comparison breaks
this confound. scVI is a foundation model (variational autoencoder
trained on millions of cells) that outputs $d = 50$ embeddings; it
clears the null in 16 of 25 tissues. Bag-of-genes PCA at $d = 512$---a
classical method, not a foundation model---fails in all 25
(Table~\ref{tab:confound}). Both model families fail at $d = 512$; the
foundation model succeeds at $d = 50$. Dimensionality is the causal
variable.

\begin{table}[t]
\centering
\caption{Confound control: model family vs.\ dimensionality. Fraction of
25 tissues where CKA exceeds the analytic null. Rows: model family.
Columns: embedding dimension.}
\label{tab:confound}
\small
\begin{tabular}{@{}lcc@{}}
\toprule
& $d = 50$ & $d = 512$ \\
\midrule
Foundation model & scVI: 16/25 (64\%) & Geneformer: 0/25 (0\%) \\
PCA / random     & (external: 86\%)   & BoG-PCA: 0/25 (0\%) \\
\bottomrule
\end{tabular}
\end{table}

A PCA-reduction control confirms this. Geneformer ($d = 512$) and
scGPT ($d = 512$) cell-level embeddings were PCA-reduced to $d = 50$
and $d = 200$ (joint PCA on source and target cells), and CKA was
recomputed against the $(k, d_{\mathrm{reduced}})$ null. At native
$d = 512$, Geneformer clears the null in 0/22 tissues and scGPT in
0/22. After reduction to $d = 50$, Geneformer clears 18/22 (82\%) and
scGPT 8/22 (36\%). At $d = 200$, Geneformer clears only 2/22 and scGPT
0/22 (Table~\ref{tab:pca_reduction}). The same embeddings, from the
same model, cross the certification boundary when their dimensionality
crosses it. Model architecture is held constant; dimensionality is the
causal variable.

\begin{table}[t]
\centering
\caption{PCA-reduction confound control. Fraction of 22 tissues where
CKA exceeds the $(k, d)$-matched analytic null, by model and
dimensionality. scVI ($d = 50$) shown for reference (no reduction
applied).}
\label{tab:pca_reduction}
\small
\begin{tabular}{@{}lccc@{}}
\toprule
Model & Native $d$ & $d = 50$ & $d = 200$ \\
\midrule
Geneformer      & 0/22 (0\%)  & 18/22 (82\%) & 2/22 (9\%) \\
scGPT           & 0/22 (0\%)  & 8/22 (36\%)  & 0/22 (0\%) \\
scVI (ref.)     & 15/22 (68\%) & ---         & --- \\
\bottomrule
\end{tabular}
\end{table}

Null-normalizing CKA by subtracting the analytic mean and dividing by
the null standard deviation weakens rather than improves predictive
validity ($z$-CKA: $\rho = +0.42$ vs.\ raw CKA $\rho = +0.46$;
partial $\rho|d$: $0.39$ vs.\ $0.44$). Within the $d = 50$ stratum,
the degradation is larger ($\rho = +0.17$ vs.\ $+0.45$): $z$-scoring
removes the $k$-dependent signal that CKA carries at low $d$, where the
metric has real discriminative power. The normalization is actively
harmful in the regime where CKA works and redundant in the regime where
it does not. The constructive fix is the usability criterion itself,
not a rescaled score.


\subsection{The null floor is structural: pre-registered spectral-gap test}
\label{sec:spectral}

The null floor could reflect either a normalization artifact fixable by
a better statistic, or a structural property of small-$k$ centroid
comparison. To distinguish these, we pre-registered a spectral-gap
ratio (SGR) designed to exploit eigenvalue structure rather than
Gram-matrix alignment (pre-registration SHA-256 frozen before data
access; Supplementary~\ref{app:spectral}).

Random $k \times d$ matrices have flat singular-value spectra
(Marchenko--Pastur distribution). Cell-type centroid matrices from
trained models should have concentrated spectra---a few dominant
directions capturing biological axes such as immune versus
epithelial---producing higher explained-variance ratios than random.
For a $k \times d$ centroid matrix $C$, the spectral-gap ratio is
$\mathrm{SGR}_1 = s_1^2 / \sum_i s_i^2$ (fraction of variance in the
top singular value of the centered matrix). The discriminative
statistic is a $z$-score against $k,d$-matched null distributions
(5,000 Gaussian draws per pair).

The experiment produces a clean negative result across 22 tissues,
136 conditions, and 8 embedding models. Null baselines (random
projection, untrained encoder) exhibit \emph{higher} spectral
concentration than trained models: baseline mean $z_1 = 141$,
contender mean $z_1 = 88$ (Mann--Whitney $U = 1{,}058$, one-sided
$p > 0.99$; Table~\ref{tab:sgr}). SGR does not predict transfer
(partial $\rho = -0.11$, $p = 0.99$) and adds no information beyond
CKA (partial $\rho = 0.001$, $p = 0.99$).

\begin{table}[t]
\centering
\caption{Pre-registered spectral-gap test results (SGR$_1$ $z$-scores
against $k,d$-matched Gaussian null). Higher $z$ = more concentrated
spectrum. Trained models have \emph{lower} spectral concentration than
baselines, opposite to prediction.}
\label{tab:sgr}
\small
\begin{tabular}{@{}lcccl@{}}
\toprule
Model & $n$ & $z_1$ mean & SGR$_1$ mean & Role \\
\midrule
bog\_pca\_512           & 22 & 167 & 0.529 & contender \\
untrained encoder       & 22 & 155 & 0.498 & baseline \\
random projection       & 22 & 126 & 0.434 & baseline \\
scGPT                   & 22 & 76  & 0.311 & contender \\
Geneformer              & 22 & 65  & 0.275 & contender \\
scVI                    & 22 & 18  & 0.364 & contender \\
\bottomrule
\end{tabular}
\end{table}

The reversal has a clear mechanism. Random projections and untrained
encoders preserve the dominant variance axis of the input gene-expression
matrix (the immune--epithelial split that dominates transcriptomic
variation). Trained models redistribute variance across embedding
dimensions to separate cell types more evenly---precisely the
representation geometry that improves downstream classification but
lowers spectral concentration. The null floor is structural: it
arises because \emph{any} projection of biologically structured
gene-expression data produces centroid matrices with concentrated
spectra, regardless of whether the model learned meaningful
representations.


\subsection{Bio-conservation metrics under controlled degradation}
\label{sec:noise}

A metric suitable for certification should decrease monotonically as the
embedding is degraded. We tested this by adding Gaussian noise at seven
levels ($\sigma \in \{0.01, 0.05, 0.1, 0.25, 0.5, 1.0, 2.0\} \times$
embedding standard deviation) to six embedding methods across four
tissues (24 conditions, single noise seed per condition).

Every bio-conservation metric shows non-monotonic trajectories in this
single-seed experiment (Table~\ref{tab:noise}). ARI is non-monotonic in
22 of 24 conditions; NMI in 17; graph connectivity in 22; cLISI in 11.
At intermediate noise levels, scores \emph{increase} before collapsing
at high noise. The mechanism differs by metric: noise drives cells into
fewer Leiden clusters, mechanically raising NMI; noise smooths local
neighborhoods, transiently increasing graph connectivity. Both
batch-correction metrics with computable values (silhouette-batch, PCR
comparison) pass the monotonicity test (23/24 each), but for a
mechanical reason: noise destroys batch structure, so batch-mixing
scores improve regardless of embedding quality.

\begin{table}[t]
\centering
\caption{Noise dose-response (single seed). Non-monotonicity count
across 24 conditions (4 tissues $\times$ 6 embeddings) at 7 noise
levels. A metric passes if $\leq 1$ condition is non-monotonic. ARI
and NMI depend on Leiden clustering, which is sensitive to
perturbation; multi-seed replication is needed to confirm that
non-monotonicity persists in expectation.}
\label{tab:noise}
\small
\begin{tabular}{@{}lccc@{}}
\toprule
Metric & Non-monotonic & Total & Verdict \\
\midrule
ARI (Leiden)         & 22 & 24 & FAIL \\
Graph connectivity   & 22 & 24 & FAIL \\
NMI (Leiden)         & 17 & 24 & FAIL \\
cLISI                & 11 & 24 & FAIL \\
Isolated label ASW   &  9 & 24 & FAIL \\
Silhouette (label)   &  3 & 24 & FAIL \\
iLISI                &  2 & 24 & FAIL \\
\midrule
Silhouette (batch)   &  1 & 24 & PASS \\
PCR comparison       &  1 & 24 & PASS \\
\bottomrule
\end{tabular}
\end{table}

Non-monotonicity, if confirmed across seeds, would mean that a metric
can report improvement when an embedding is degraded---a third instance
of the ranking/certification dissociation. Combined with the analytic
null-floor result, the pattern suggests that metrics built from pairwise
distances, neighbor graphs, and cluster geometry respond to geometric
properties of the embedding space that do not track biological content
monotonically. This result is preliminary pending multi-seed
replication.


\subsection{Cross-tissue replication localizes the failure}
\label{sec:crosstissue}

The preceding results all evaluate metrics within the cross-assay
setting: source and target come from the same tissue but different
sequencing protocols (10x Chromium vs.\ Smart-seq2). Protocol
differences introduce technical artifacts---library-size distributions,
gene-detection rates, cell-type proportions---that could confound metric
behavior. To test whether the ranking failures are specific to
cross-assay evaluation, we evaluated scIB bio-conservation metrics on
\emph{cross-tissue} transfer: training a classifier on one tissue's
embeddings and testing on another (5 tissues, 6 tissue pairs, 4 models,
24 conditions).

scIB metrics predict cross-tissue transfer with substantially lower
inversion rates than cross-assay transfer. The grand-mean pairwise
inversion rate (fraction of model pairs where the metric ranks
$A > B$ but transfer F1 ranks $B > A$) is 22.1\% across five scIB
metrics (bootstrap 95\% CI: 13.1--34.9\%;
Table~\ref{tab:crosstissue}). The corresponding cross-assay inversion
rates from the same metrics on the main panel range from 30\% to 48\%
(Table~\ref{tab:crossassay_inv}). Spearman correlation between
aggregate scIB scores and transfer F1 is $\rho = +0.63$ ($p = 0.0014$,
permutation test; $n = 24$ conditions).

\begin{table}[t]
\centering
\caption{Cross-tissue pairwise inversion rates for scIB
bio-conservation metrics (6 tissue pairs, 4 models, 35 pairwise
comparisons per metric). Lower is better; 50\% = coin flip. Block
bootstrap 95\% CIs resample tissue pairs.}
\label{tab:crosstissue}
\small
\begin{tabular}{@{}lccc@{}}
\toprule
Metric & Inversion rate & 95\% CI & $n$ pairs \\
\midrule
Silhouette (label)   & 17.1\% & [6.1, 27.8] & 35 \\
Isolated label ASW   & 17.1\% & [8.9, 25.0] & 35 \\
NMI (Leiden)         & 22.9\% & [11.1, 36.1] & 35 \\
ARI (Leiden)         & 25.7\% & [8.9, 50.0] & 35 \\
cLISI                & 28.6\% & [13.9, 38.9] & 35 \\
\midrule
Grand mean           & 22.1\% & [13.1, 34.9] & --- \\
\bottomrule
\end{tabular}
\end{table}

\begin{table}[t]
\centering
\caption{Cross-assay inversion rates for the same scIB metrics on the
main evaluation panel (4 tissues, 6 models, 60 comparisons per metric).
Cross-assay rates are uniformly higher than cross-tissue rates
(Table~\ref{tab:crosstissue}).}
\label{tab:crossassay_inv}
\small
\begin{tabular}{@{}lcc@{}}
\toprule
Metric & Inversion rate & $n$ pairs \\
\midrule
NMI (Leiden)         & 30.0\% & 60 \\
ARI (Leiden)         & 31.7\% & 60 \\
cLISI                & 35.0\% & 60 \\
Silhouette (label)   & 43.3\% & 60 \\
Isolated label ASW   & 48.3\% & 60 \\
\bottomrule
\end{tabular}
\end{table}

The result is seed-stable. Three independent cell-sampling seeds produce
inversion rates of 22.1\%, 22.2\%, and 17.8\%, and Spearman
correlations of $\rho = +0.63$, $+0.76$, and $+0.74$ (all
$p < 0.002$). The cross-tissue advantage is robust to the specific
cells sampled.

The contrast localizes the source of metric unreliability. In the
cross-assay setting, protocol artifacts (different sequencing
technologies applied to the same tissue) confound scIB scores: metrics
that respond to library-size distributions and gene-detection rates
can disagree with downstream transfer. In the cross-tissue setting, both
source and target use the same sequencing protocol; the variation is
biological (different tissues), and scIB metrics track this variation
more faithfully. The ranking failure is protocol-dependent, not
intrinsic to the metrics.


\subsection{External validation}
\label{sec:external}

Cell-type CKA was tested on three datasets with no overlap with the
Census evaluation: pancreas \citep{luecken2022scib} (6 technologies, 56
assay pairs, $k = 12$--$14$, $d = 50$ and $200$), Tabula Sapiens
\citep{tabulasapiens2022} (10x and Smart-seq2 across 8 organs, 16
pairs, $k = 9$--$22$, $d = 50$ and $200$), and PBMC
\citep{ding2020systematic} (7 technologies, 28 pairs, $k = 8$--$9$,
$d = 50$ and $200$). Embedding models are PCA-only baselines (no
foundation models), isolating metric behavior from model choice. These
datasets operate at lower dimensionality than the Census panel,
providing the external arm of the discriminability gradient
(Table~\ref{tab:deffect}).

CKA generalizes to two of three datasets
(Table~\ref{tab:external}). In pancreas, CKA achieves partial
$\rho = +0.59$ ($p = 0.0001$); in Tabula Sapiens, $\rho = +0.60$
($p = 0.035$). Procrustes similarity is consistently positive across all
three ($\rho = +0.56$, $+0.73$, $+0.11$), matching the Census direction
but reaching significance in none---a power limitation from fewer
technology pairs per dataset. PBMC shows a ceiling effect: all seven
technologies produce CKA $> 0.92$ on canonical cell types (T cells, B
cells, monocytes), leaving insufficient variance for discriminative
ranking ($\rho = +0.09$, $p = 0.87$; Supplementary~\ref{app:pbmc}). CKA
predicts transfer when cell-type difficulty varies across technologies
and fails when the task is uniformly easy.

\begin{table}[t]
\centering
\caption{External validation on three independent datasets (no Census
data, no foundation models). Partial Spearman $\rho$ controlling for
embedding dimensionality; tissue-stratified permutation $p$-values.
Bold: $p < 0.05$.}
\label{tab:external}
\small
\begin{tabular}{@{}lcccccc@{}}
\toprule
& \multicolumn{2}{c}{Pancreas} & \multicolumn{2}{c}{Tabula Sapiens} & \multicolumn{2}{c}{PBMC} \\
\cmidrule(lr){2-3}\cmidrule(lr){4-5}\cmidrule(lr){6-7}
Metric & $\rho$ & $p$ & $\rho$ & $p$ & $\rho$ & $p$ \\
\midrule
Cell-type CKA & $+0.592$ & $\mathbf{0.0001}$ & $+0.604$ & $\mathbf{0.035}$ & $+0.089$ & 0.874 \\
Procrustes    & $+0.563$ & 0.115              & $+0.726$ & 0.218             & $+0.113$ & 0.659 \\
\midrule
Silhouette (tgt) & $+0.172$ & $\mathbf{<0.001}$ & $+0.664$ & $\mathbf{0.032}$ & $+0.627$ & $\mathbf{<0.001}$ \\
MMD              & $-0.451$ & 1.000              & $-0.275$ & 0.744             & $-0.135$ & 1.000 \\
Domain AUC       & $-0.495$ & 0.095              & $+0.015$ & 0.839             & $-0.148$ & 0.801 \\
\bottomrule
\end{tabular}
\end{table}


\section{Discussion}
\label{sec:discussion}

Ranking and certification are distinct properties of embedding metrics,
and the central finding is a computable boundary between regimes where
CKA can and cannot certify. Given $k$ shared cell types and embedding
dimension $d$, the expected CKA between random Gaussian centroid
matrices is $\mathbb{E}[\mathrm{CKA}] \approx 1 -
1.06\,(k{-}1)/(d{+}k)$. CKA can certify biological signal only when
observed scores exceed this value. At $d = 512$, the threshold requires
roughly 168 shared cell types, beyond any existing single-cell dataset.
The failure is arithmetic rather than empirical.

Near-orthogonality of high-dimensional random vectors is textbook
\citep{johnsonlindenstrauss1984}; the non-obvious part is where the
boundary lands. The crossover from certifiable ($d = 50$, 86\% of
conditions clear the null) to uncertifiable ($d \geq 512$, 0\% clear)
happens inside the $d = 50$--$512$ range that the field currently uses.
PCA-based methods ($d = 50$--$200$) were on the correct side; foundation
models ($d \geq 512$) crossed to the wrong side. The PCA-reduction
control confirms this is causal: Geneformer embeddings PCA-reduced from
$d = 512$ to $d = 50$ clear the null in 18/22 tissues, versus 0/22 at
native dimensionality. A researcher using CKA to evaluate scVI
($d = 50$) gets real certification signal; the same researcher using the
same metric on Geneformer ($d = 512$) gets none, and the metric gives no
warning that it has failed.

The pre-registered spectral-gap experiment confirms that the null floor
is structural. The spectral-gap ratio was designed to exploit eigenvalue
structure to distinguish trained from random embeddings. Trained models
show \emph{less} spectral concentration than random projections, the
opposite of the pre-registered prediction. The reversal arises because
trained models spread cell types across embedding dimensions to improve
separability, while random projections preserve the input's dominant
variance axis. Any projection of biologically structured gene expression
produces concentrated centroid spectra; the concentration reflects input
structure, not learned representations.

Null-normalizing CKA by $z$-scoring against the analytic null is
actively counterproductive. At $d = 50$, where CKA has real
discriminative power, the normalization removes the $k$-dependent signal
that drives the correlation with transfer (within-stratum $\rho$ drops
from $+0.45$ to $+0.17$). At $d = 512$, where CKA is uninformative,
the normalization is redundant. The fix is the criterion itself---a
researcher who checks whether their $(k, d)$ pair permits discrimination
before computing CKA avoids the problem entirely.

The same geometric confound appears in neural population geometry, where
19 metrics are confounded by neuron count \citep{tower2026bracket}. The
recurrence across biological domains suggests that geometric evaluation
in high dimensions requires null-model calibration whenever $d \gg k$.

The noise dose-response provides preliminary evidence of a
complementary certification failure. Bio-conservation metrics respond
non-monotonically to controlled degradation in single-seed
experiments---scores increase at intermediate noise before
collapsing. Because ARI and NMI depend on Leiden clustering, which is
sensitive to small perturbations, multi-seed replication is needed to
confirm that the non-monotonicity persists in expectation rather than
reflecting clustering jitter.

The cross-tissue replication experiment identifies where scIB metrics
recover. Pairwise inversion rates drop from 30--48\% in the cross-assay
setting to 22\% in the cross-tissue setting ($\rho = +0.63$,
$p = 0.0014$, stable across three seeds). The cross-assay setting
confounds metric scores with protocol artifacts: library-size
distributions, gene-detection rates, and cell-type proportions differ
systematically between 10x Chromium and Smart-seq2. The cross-tissue
setting removes this confound---source and target use the same
sequencing protocol, and the variation is biological. scIB metrics track
biological variation faithfully when protocol artifacts are absent. The
ranking failure is a property of the evaluation setting, not of the
metrics themselves.

Procrustes similarity achieves the strongest transfer prediction in the
Census panel ($\rho = +0.66$) and is consistently positive across all
three external datasets ($\rho = +0.56$, $+0.73$, $+0.11$). CKA
achieves significance in two of three external datasets. Target-domain
silhouette is significant in all three external datasets but is likely
confounded: target assays with well-separated cell types are easier to
classify regardless of source embedding quality.

\paragraph{Practical recommendations.} First, before interpreting CKA as
evidence of biological signal, compute $\mathbb{E}[\mathrm{CKA}]$ at
your $k$ and $d$ using Equation~\ref{eq:cka_null}; if observed scores do
not exceed this value, the metric can rank but cannot certify. Second,
use CKA and Procrustes for ranking across models, controlling for
embedding dimensionality; do not interpret absolute scores as indicating
the presence or absence of biological signal at $d \geq 512$. Third,
always include matched-dimensionality null baselines (random projection,
untrained encoder) in evaluation panels. Fourth, test metric behavior
under controlled degradation before adopting a new evaluation metric.
Fifth, prefer Procrustes similarity over CKA for cross-technology
transfer prediction when cell-type labels are available.

\paragraph{Limitations.} All evaluations use cross-technology cell-type
transfer F1 as ground truth; alternative criteria such as perturbation
prediction may reveal different metric properties. The Census panel
spans 25 tissues and eight embedding methods across three architectural
families (transformer, VAE, PCA); expansion to additional models would
increase statistical power. The PCA-reduction control confirms the causal role of dimensionality
for Geneformer (18/22 tissues clear the null at $d = 50$), though
scGPT's lower recovery rate (8/22) suggests that model architecture
modulates the effect size within the certifiable regime. The noise
dose-response experiment
uses four tissues, six methods (24 conditions), and a single noise seed
per condition; because ARI and NMI depend on Leiden clustering, which is
sensitive to perturbation, multi-seed replication is needed to
distinguish genuine non-monotonicity from clustering jitter. The
cross-tissue experiment uses five tissues and four models (24
conditions); expanding to more tissues and models would tighten the
inversion-rate estimates. The PBMC ceiling effect limits
external validation to two of three datasets. The spectral-gap test is
one of many possible statistics; other approaches might rescue
discriminative validity. The closed-form CKA null is exact for Gaussian
inputs; real centroid distributions are non-Gaussian, though the
simulation match suggests the approximation is tight in practice. We
release all scripts, pre-registration records, and result files to
support replication.


\section{Methods}

\subsection{Data}

All embeddings are accessed via the CellxGene Census API (v2023-12-15)
\citep{abdulla2025cellxgene}. Cross-technology pairs compare 10x
Chromium 3$'$ v3 (EFO:0009922, source) against Smart-seq2
(EFO:0008931, target) within the same tissue, with zero donor overlap
enforced by metadata filtering. Tissues were auto-discovered by
requiring $\geq 200$ cells per assay and $\geq 8$ shared cell types,
yielding 25 tissues. All evaluations subsample $n = 2{,}000$ cells per
side.

Eight embedding models: Geneformer ($d = 512$)
\citep{theodoris2023geneformer}, Geneformer-v2 104M and 316M,
scGPT ($d = 512$) \citep{cui2024scgpt}, scVI ($d = 50$)
\citep{lopez2018scvi}, bag-of-genes PCA-512 ($d = 512$;
$\log(1 + x)$ without total-count normalization, PCA with fixed seed;
following \citealp{cole2026morpheus}), random Gaussian projection
($R \in \mathbb{R}^{512 \times G}$, $R_{ij} \sim \mathcal{N}(0,
1/\sqrt{512})$), and an untrained two-layer MLP encoder (random weights,
$d = 512$). Six contender models (excluding random projection and
untrained encoder) produce 104 contender conditions.

\paragraph{External validation datasets.} Three independent datasets
were used for external validation (Section~\ref{sec:external}), each
with PCA-only embeddings. Pancreas: the scIB benchmark dataset
\citep{luecken2022scib}, six technologies, 56 assay pairs, 224
contender conditions. Tabula Sapiens \citep{tabulasapiens2022}: 10x and
Smart-seq2 across eight organs, 16 pairs, 64 contenders. PBMC
\citep{ding2020systematic}: seven technologies, 28 pairs, 112
contenders.

\subsection{Cell-type centroid metrics}

Cell-type CKA \citep{kornblith2019similarity} computes centered kernel
alignment on $k \times d$ centroid matrices (mean embedding per shared
cell type) between source and target assays, using linear kernels.
Procrustes similarity applies orthogonal Procrustes analysis to
column-centered centroid matrices after PCA reduction to $\min(50,
k-1, d)$ dimensions, and returns $1 - d_{\mathrm{Proc}}^2$ after
optimal rotation. Both metrics require cell-type labels and are
evaluated on the $k$ cell types shared between source and target.

\subsection{Random centroid null}

For each pair $(k, d)$ observed in the panel, we generate 1,000 pairs
of random Gaussian matrices ($k \times d$, entries $\sim
\mathcal{N}(0,1)$) and compute CKA and Procrustes similarity. Each real
condition is compared to its $(k, d)$-matched null distribution. The
analytic null (Equation~\ref{eq:cka_null}) provides a closed-form
alternative verified within 0.02\% of simulation across $k = 8$--$44$
and $d = 50$--$1{,}280$. Null-normalized scores are computed as
$(x - \mu_{k,d}) / \sigma_{k,d}$ where $\mu_{k,d}$ and $\sigma_{k,d}$
are the null mean and standard deviation at the condition's $(k, d)$.

\subsection{Closed-form CKA null}
\label{sec:closedform}

For two independent random Gaussian matrices $X, Y \in \R^{k \times d}$,
linear CKA is $\mathrm{CKA} = \|H_k X Y^\top H_k\|_F^2 / (\|H_k X
X^\top H_k\|_F^2 \cdot \|H_k Y Y^\top H_k\|_F^2)^{1/2}$ where
$H_k = I_k - \frac{1}{k}\mathbf{1}\mathbf{1}^\top$. The denominator
factors into two identical terms by symmetry, so
$\mathbb{E}[\mathrm{CKA}] \approx \mathbb{E}[\mathrm{tr}((HXY^\top
H)^2)] / \mathbb{E}[\mathrm{tr}((HXX^\top H)^2)]$.

Since $W = XX^\top$ follows a Wishart distribution
$\mathcal{W}_k(d, I_k)$, the moments of the centered Gram matrix $HWH$
are computable from Wishart moment identities. We use
$\mathbb{E}[\mathrm{tr}(W^2)] = kd(1 + k + d)$, verified empirically
within 0.2\%. Centering reduces the effective rank by one degree of
freedom. Fitting across 15 values of $k$ and $d = 512$ yields the
approximation in Equation~\ref{eq:cka_null} with coefficient of
variation 2.0\% on the scaling constant $c$.

\subsection{Spectral-gap ratio}
\label{sec:sgrmethod}

For a $k \times d$ centroid matrix $C$ (mean source embedding per
shared cell type), center by subtracting the row mean and compute
singular values $s_1 \geq s_2 \geq \cdots \geq s_r$ with $r =
\min(k{-}1, d)$. The spectral-gap ratio at rank $j$ is $\mathrm{SGR}_j
= \sum_{i=1}^j s_i^2 / \sum_{i=1}^r s_i^2$. For each $(k, d)$ pair in
the panel, 5,000 random Gaussian $k \times d$ matrices define the null
distribution of $\mathrm{SGR}_j$ (seed 42). The $z$-score is
$({\mathrm{SGR}_j - \mu_{\mathrm{null}}}) / {\sigma_{\mathrm{null}}}$.
The null distributions were frozen (SHA-256) before any real data were
accessed.

\subsection{PCA-reduction control}

For each tissue, Geneformer and scGPT cell-level embeddings ($d = 512$)
are PCA-reduced to $d_{\mathrm{out}} \in \{50, 200\}$ by fitting PCA on
the concatenation of source and target cells (joint PCA). Cell-type CKA
is recomputed on the reduced centroids and compared to the
$(k, d_{\mathrm{out}})$-matched analytic null. scVI embeddings
($d = 50$) are included without reduction as a reference.

\subsection{Cross-tissue transfer}
\label{sec:crosstissuemethod}

Five tissues (lung, liver, kidney, brain, blood) were selected from the
Census panel, yielding six directional tissue pairs (lung$\to$brain,
lung$\to$kidney, blood$\to$lung, blood$\to$brain, liver$\to$kidney,
blood$\to$liver). Four embedding models (Geneformer, scVI, scGPT,
bag-of-genes PCA-512) produce 24 conditions. For each condition, scIB
bio-conservation metrics are computed on the combined source and target
embeddings. Transfer F1 (macro-averaged logistic regression on shared
cell types) is the ground truth. Pairwise inversion rates count the
fraction of model pairs where a metric disagrees with F1 on which model
is better, computed per tissue pair. The grand-mean inversion rate
averages across metrics and tissue pairs; block bootstrap (10,000
resamples of tissue pairs) provides 95\% CIs. Three independent
cell-sampling seeds (20260801, 20260802, 20260803) test stability.

\subsection{Statistical methods}

All correlations are Spearman rank with $10{,}000$-resample bootstrap
95\% CIs. Partial Spearman correlations control for embedding
dimensionality via rank residuals from OLS on ranked variables.
Tissue-stratified permutation tests (10,000 iterations) permute F1
values within tissues independently and report the proportion of
permutations with $|\rho_{\mathrm{perm}}| \geq |\rho_{\mathrm{obs}}|$.

\subsection{Pre-registration}

Scorer source code, hyperparameters, dataset specifications, and
hypothesis thresholds were frozen via SHA-256 hashing before data
access. SHA-256 hashes and the full deviation log are in the online
repository.


\section*{Data Availability}

Scorer source code, pre-registration records (SHA-256 hashes),
experiment scripts, and raw result JSON files are available at
\url{https://github.com/elliottower/preflight-bio}. All embedding data
are accessed via the CellxGene Census API.

\section*{Code Availability}

All analysis scripts are available at the repository listed under Data
Availability.

\section*{Declarations}

\paragraph{Competing interests.} The author declares no competing
interests.

\paragraph{Funding.} None. This work was conducted independently.

\paragraph{Authors' contributions.} E.T.\ conceived the study,
designed and pre-registered the protocol, performed all analyses, and
wrote the manuscript.

\paragraph{Acknowledgements.} AI tools (Claude Code) were used for code
drafting and manuscript drafting. All experimental design and conclusions
are the author's.


\appendix

\section{PBMC ceiling effect}
\label{app:pbmc}

The PBMC dataset \citep{ding2020systematic} provides seven scRNA-seq
technologies applied to the same donor PBMCs. Cell-type CKA ranges from
0.92 to 0.97 across all 28 source-target technology pairs. This
variance compression explains the null correlation with transfer F1
($\rho = +0.09$, $p = 0.87$): the metric cannot rank technologies on
a property where they score nearly identically. T cells, B cells, NK
cells, and monocytes are well-separated in any reasonable embedding; no
technology fails to capture this structure. The pancreas dataset contains
rare endocrine subtypes that are harder to resolve, providing sufficient
CKA variance for the correlation to emerge.


\section{Geometric modules}
\label{app:modules}

Five geometric modules for evaluating embedding transportability were
pre-registered and evaluated alongside the main metrics. Four were
eliminated by construct-validity checks: M1 (subspace overlap) and M6
(graph curvature) are confounded by embedding dimensionality via
Johnson--Lindenstrauss distance preservation; M2 (domain separability)
anti-predicts transfer; M4 (participation ratio) has
dimension-dependent normalization artifacts. M3 (direction stability)
was the sole candidate to survive construct-validity screening, but
achieves $\rho = +0.02$ ($p = 0.83$) on the full panel---no
predictive validity. These results demonstrate that the validity
protocol is non-trivial: it eliminates metrics, including those designed
by the protocol's authors.


\section{Pre-registered spectral-gap experiment}
\label{app:spectral}

The spectral-gap experiment was pre-registered with four hypotheses and
explicit decision rules for all pass/fail combinations. The full
pre-registration document, SHA-256 hashes of frozen null distributions
and analysis scripts, and all raw results are in the online repository.

\paragraph{Decision rules (from pre-registration).} HS1 (discriminative):
contender models have higher SGR $z$-scores than baselines (Mann--Whitney
$p < 0.05$). HS2 (predictive): SGR $z$-score correlates with transfer
F1 (partial $\rho > 0.20$, permutation $p < 0.05$). HS3 (incremental):
SGR adds variance beyond CKA. HS4 (both axes): HS1 and HS2
simultaneously.

\paragraph{Results.} HS1 FAIL: baselines have higher $z$-scores
(mean 141 vs.\ 88; $p > 0.99$). HS2 FAIL: partial $\rho = -0.11$
($p = 0.99$). HS3: partial $\rho = 0.001$ ($p = 0.99$). HS4 FAIL by
inheritance.

\paragraph{Interpretation (pre-registered ``both fail'' rule).} Spectral
concentration does not solve the discriminative-validity problem. The
failure is informative: it eliminates eigenvalue normalization as a fix
and establishes that the deficiency is structural to small-$k$ centroid
inputs. Any projection of biologically structured data produces
centroid matrices with concentrated spectra, regardless of whether the
projection learned meaningful representations.


\section{Noise dose-response: full data}
\label{app:noise}

Table~\ref{tab:app_noise_fail} shows representative non-monotonic
trajectories. NMI in brain/Geneformer shows scores \emph{increasing} at
high noise ($\sigma \geq 1.0$): noise drives all cells into fewer
Leiden clusters, mechanically raising NMI.

\begin{table}[h!]
\centering
\caption{Representative non-monotonic trajectories. Values are metric
scores at $\sigma \in \{0.01, 0.05, 0.1, 0.25, 0.5, 1.0, 2.0\} \times
\text{embedding\_std}$.}
\label{tab:app_noise_fail}
\footnotesize
\begin{tabular}{@{}llp{0.6\textwidth}@{}}
\toprule
Metric & Condition & $\sigma$: 0.01, 0.05, 0.1, 0.25, 0.5, 1.0, 2.0 \\
\midrule
Graph conn. & brain / scgpt & 0.772, 0.771, 0.773, \textbf{0.777}, \textbf{0.793}, 0.748, 0.589 \\
NMI & brain / geneformer & 0.550, 0.550, 0.545, \textbf{0.547}, 0.545, \textbf{0.575}, \textbf{0.577} \\
ARI & kidney / scvi & 0.397, \textbf{0.410}, 0.394, 0.351, 0.224, 0.073, 0.012 \\
cLISI & lung / scvi & 0.999, 0.999, \textbf{0.999}, 0.998, 0.975, 0.871, 0.642 \\
\bottomrule
\end{tabular}
\end{table}


\bibliographystyle{plainnat}
\bibliography{paper_d_v11}

\end{document}
